\chapter{Prove the Database Boundary from Another Jail}

The database health check proves that PostgreSQL is alive inside its own jail.
It does not prove that another service can route to it, reach port 5432, resolve
package repositories, authenticate, or complete a query. This chapter adds the
future catalog service as a blank client jail and checks those layers one at a
time.

The catalog jail is deliberately not a Rails application yet. Keeping it small
lets us diagnose the network and database boundary before Ruby, Bundler, Rails,
and application configuration add more possible causes of failure.

\section{The proposed client}

Use these values for the second jail:

\begin{description}
  \item[starter] Blank;
  \item[jail name] \texttt{catalog};
  \item[FreeBSD release] \texttt{15.1-RELEASE};
  \item[private address] \texttt{10.17.89.24/24};
  \item[outbound policy] Internet only;
  \item[starts at boot] disabled; and
  \item[inbound exposure] none.
\end{description}

Internet-only access is needed temporarily to install the PostgreSQL client and,
in the next chapter, the Ruby toolchain. It does not create an inbound port
forward. The catalog and database jails share BSDRunner's managed Blackship
network, so the catalog can address PostgreSQL directly at
\texttt{10.17.89.23:5432}.

\section{Exercise 4A: Create the blank catalog jail}

Open BSDRunner Jails and choose \textbf{Create a Service Jail}, then select the
\textbf{Blank} starter. Review the same five concepts used for the database
jail:

\begin{enumerate}
  \item reuse the existing Blackship project and ZFS storage;
  \item select \texttt{15.1-RELEASE};
  \item keep the proposed address \texttt{10.17.89.24/24};
  \item name the jail \texttt{catalog}, select \textbf{Internet only}, and
    leave \textbf{Start at boot} disabled; and
  \item confirm that the plan has no inbound port forward before creating it.
\end{enumerate}

If the jail already exists from an earlier session, do not create a duplicate.
Open its detail view, confirm the same address and policy, start it if necessary,
and open its console.

\section{Exercise 4B: Test only the PostgreSQL port}

From the catalog jail's console, make a bounded TCP connection attempt:

\begin{lstlisting}[language=sh]
nc -vz -w 3 10.17.89.23 5432
\end{lstlisting}

The expected result ends with:

\begin{lstlisting}
Connection to 10.17.89.23 5432 port [tcp/postgresql] succeeded!
\end{lstlisting}

This result proves three things: the catalog jail has a route to the database
jail, PostgreSQL is listening on its private address, and no intervening filter
rejects the connection. It does not test a database name, user, password, or
SQL query.

Stop here if the connection times out or is refused. Check that both jails are
running and still have the addresses assigned in this tutorial before changing
PostgreSQL authentication or installing application software.

\section{Exercise 4C: Install a database client}

Still inside the catalog jail, verify DNS before asking the package manager to
bootstrap or download anything:

\begin{lstlisting}[language=sh]
drill pkg.FreeBSD.org
\end{lstlisting}

The response should contain \texttt{rcode: NOERROR} and identify the upstream
resolver selected by BSDRunner. In this lab that resolver is
\texttt{192.168.1.254}; it is not a DNS service running inside the jail.

Install only the PostgreSQL client tools:

\begin{lstlisting}[language=sh]
pkg install -y postgresql18-client
\end{lstlisting}

This is an exploratory installation. Once the catalog application shape is
known, its packages and gems will move into a repeatable build definition.

If package bootstrap reports a resolver failure, do not replace the host's DNS
configuration. Inspect the jail first:

\begin{lstlisting}[language=sh]
cat /etc/resolv.conf
route -n get default
\end{lstlisting}

The resolver should match the one selected by BSDRunner and the default gateway
should be \texttt{10.17.89.1}. On the host, the read-only command
\texttt{ifconfig blackship0 inet} should show that gateway with the
\texttt{/24} prefix. A guarded jail start runs \texttt{blackship setup}, which
reconciles an existing bridge as well as creating a missing one.

\section{Exercise 4D: Authenticate and identify both endpoints}

Connect from the catalog jail with the temporary starter credential captured
in Chapter 3:

\begin{lstlisting}[language=sh]
psql -h 10.17.89.23 -U app -d appdb
\end{lstlisting}

Enter the password at the prompt. Do not put it on the command line or in this
book. At the \texttt{appdb=>} prompt, run:

\begin{lstlisting}[language=SQL]
SELECT
  current_database(),
  current_user,
  inet_client_addr(),
  inet_server_addr();
\end{lstlisting}

Verify the four returned values:

\begin{description}
  \item[database] \texttt{appdb};
  \item[user] \texttt{app};
  \item[client address] \texttt{10.17.89.24}; and
  \item[server address] \texttt{10.17.89.23}.
\end{description}

Exit with \texttt{\textbackslash q}. The endpoint addresses are important:
they demonstrate that the authenticated session crossed directly from the
catalog jail to PostgreSQL over the managed private network rather than through
an inbound host port forward.

\section{Checkpoint}

The database boundary is ready for application work when:

\begin{itemize}
  \item the catalog jail is running at \texttt{10.17.89.24};
  \item neither catalog nor database has an inbound port forward;
  \item a bounded TCP probe reaches database port 5432;
  \item the catalog jail can resolve and download FreeBSD packages; and
  \item PostgreSQL reports the expected database, user, client address, and
    server address.
\end{itemize}

These checks prove the infrastructure boundary, not the application. The next
chapter will install a Ruby toolchain and turn the blank catalog jail into a
minimal Rails API while keeping its database ownership explicit.
